\begin{table}[t]
\caption{H2 Diff-in-Disc — Full FE + DID Robustness (Probabilistic Time)} 
\fontsize{12.0pt}{14.0pt}\selectfont
\begin{tabular*}{\linewidth}{@{\extracolsep{\fill}}lccc}
\toprule
  & (1)† & (2)† & (3)† \\ 
\midrule\addlinespace[2.5pt]
Invasion & -0.124*** & -0.217*** & -0.178*** \\ 
 & (0.016) & (0.017) & (0.009) \\ 
State Owned x Treatment & 0.078+ & 0.138* & 0.074* \\ 
{} & {(0.033)} & {(0.038)} & {(0.027)} \\ 
State FEs & X & X & X \\ 
Source FEs & X & X & X \\ 
Topic FEs & X & X & X \\ 
Clusters & 60 & 42 & 60 \\ 
Num.Obs. & 27315 & 20604 & 37622 \\ 
R2 & 0.064 & 0.092 & 0.057 \\ 
R2 Adj. & 0.062 & 0.090 & 0.057 \\ 
\bottomrule
\end{tabular*}
\begin{minipage}{\linewidth}
\vspace{.05em}
+ p < 0.1, * p < 0.05, ** p < 0.01, *** p < 0.001\\
All standard errors are clustered at the source level. Model 1 uses non-military
subset with full FEs. Model 2 uses Russian-related articles with full FEs.
Model 3 is a DID specification within the optimal bandwidth. † Standard errors
computed via wild cluster restricted (WCR) bootstrap (Cameron, Gelbach \& Miller
2008) with Webb six-point weights (B=4999, null imposed), clustered by source
domain.\\
\end{minipage}
\end{table}

